% ============================================================
% chapters/08-summary.tex -- פרק 8: סיכום – מוכנים לשלב הביצוע
% Content only. All design comes from template.cls.
% Figure: figures/ch8-next-steps.tex (TikZ)
% ============================================================

\chapter{סיכום: מוכנים לשלב הביצוע}
\label{ch:summary}

\paragraph{מה הראה שלב ההכנה.} הוויכוח על תפקידה של הנצרות בנפילת רומא נמשך כבר 1,600 שנה, והדוח טען שהסיבה העיקרית לכך אינה מחסור בידע, אלא מגבלה של שיטה: מחקרים רבים בוחנים מקרה אחד, בדוגמאות שבחרו, ובמושגים שאינם מדידים (פרק~\ref{ch:problem}). סקירת הספרות חשפה פער – בספרות שסקרנו לא מצאנו השוואה של רומא המערבית, רומא המזרחית והממלכה הפרנקית באותם מדדים – והציעה כיוון: אצל הפרנקים, הכנסייה שימשה תשתית של המדינה (פרק~\ref{ch:literature}). חקר השוק הראה שהכלים הדרושים כבר קיימים בגישה פתוחה, ושחסרה רק השכבה שמחברת ביניהם (פרק~\ref{ch:market}). מבין שלוש חלופות השיטה נבחרה השיטה המשולבת – מאגר נתונים וניתוח טקסט בעזרת \en{AI}, שנבדק מול תיוג של מומחה על מדגם מבחן חתום, נפרד מקטעי הפיתוח – ובחירה זו נשארה יציבה גם כשהמשקלים השתנו (פרק~\ref{ch:alternatives}).

\paragraph{המוכנות לשלב הבא.} המחקר מוכן לביצוע: שאלת המחקר וההשערות מנוסחות כך שאפשר לאשש אותן או להפריך אותן; לכל אחת משבע הדרישות יש ערך יעד ושיטת אימות; תוכנית העבודה מחולקת לחמש חבילות, עם נתיב קריטי, שלוש אבני דרך ותקציב תכנון של כ־111 אלף ש״ח, שמתוכם כ־22 אלף ש״ח במזומן אם זמן החוקרים ניתן כתרומה אקדמית (פרק~\ref{ch:workplan}); ולארבעת הסיכונים הקריטיים נקבעה תוכנית הפחתה, שלפי אומדן התכנון תוריד את כולם מתחת לסף הקריטי – בתנאי שהפעולות יבוצעו (פרק~\ref{ch:risks}).

\paragraph{שלוש הפעולות המיידיות.} מיד לאחר אישור הדוח יבוצעו שלוש פעולות (איור~\ref{fig:next-steps}). \textbf{ראשית}, גיוס של מומחה לתקופה ושל מומחה גיבוי, וקביעת מועד התיקוף, כדי שהנתיב הקריטי לא יהיה תלוי באדם אחד (\en{R5}). \textbf{שנית}, כתיבת ספר ההגדרות ורשימת המקורות, שהם הבסיס לכל שאר העבודה. \textbf{שלישית}, מיני־פיילוט על 50 קטעים, שיעריך מוקדם ככל האפשר את הדיוק של המודל ואת העלות של הטוקנים, עוד לפני הפיילוט המלא (\en{R2}, \en{R6}).

\begin{figure}[htbp]
  \centering
  % figures/ch8-next-steps.tex -- three immediate actions after approval (chapter 8)
% Colors come from the template; fallbacks are defined here only if missing.
\providecolor{accent}{HTML}{8B2E2E}
\providecolor{accentlight}{HTML}{F4E9E1}
\begin{tikzpicture}[
  node distance=7mm,
  act/.style={draw=accent, fill=accentlight, rounded corners=2pt, text width=38mm,
              minimum height=16mm, align=center, font=\footnotesize},
  num/.style={circle, fill=accent, text=white, font=\small\bfseries, inner sep=2pt},
  arrow/.style={-{Stealth[length=2.3mm]}, thick, draw=accent}
]
  % RTL: step 1 on the right
  \node[act] (a1) {\textbf{גיוס המומחים}\\ מומחה לתקופה ומומחה גיבוי; קביעת מועד התיקוף};
  \node[act, left=of a1] (a2) {\textbf{ספר ההגדרות ורשימת המקורות}\\ משימות 1.1 ו־2.1};
  \node[act, left=of a2] (a3) {\textbf{מיני־פיילוט}\\ 50 קטעים: הערכת דיוק ועלות};
  \node[num] at (a1.north east) {1};
  \node[num] at (a2.north east) {2};
  \node[num] at (a3.north east) {3};
  \draw[arrow] (a1) -- (a2);
  \draw[arrow] (a2) -- (a3);
\end{tikzpicture}

  \caption{שלוש הפעולות המיידיות לאחר אישור הדוח}
  \label{fig:next-steps}
\end{figure}

\begin{keybox}{בחזרה לאוגוסט 410}
כשהוויזיגותים נכנסו לרומא, בני התקופה מיהרו לשאול מי אשם: האלים שננטשו, או האל החדש. המחקר המוצע לא מבקש להכריע מי צדק אז, אלא לאפשר לבדוק את השאלה באופן שיטתי – בנתונים, בשיטה שאפשר לשחזר, ובאותם מדדים לשלוש מדינות שבחרו באותה דת והגיעו לגורלות שונים כל כך.
\end{keybox}
